\subsection{Definition via the flows of the single terms}
Let $S(t)=e^{tA}$ and $S_n(t)=e^{tA_n}$. The contribution that term $n$ \emph{alone} would make to the solution is $S_n(t)u_0-u_0$. The remainder is the interaction term
\begin{equation}\label{xeq:eint-def}
\Eint^{\mathrm{add}}(t)=u(t)-u_0-\sum_{n=1}^N\big(S_n(t)u_0-u_0\big).
\end{equation}
(The superscript ``add'' records that this is the \emph{additive} definition; below we write $\Eint$ for $\Eint^{\mathrm{add}}$ unless the distinction matters.)

\subsection{Exact formula in the commuting case}
\begin{aside}{Why constant coefficients commute, and what $\Delta_n$ is}
Constant-coefficient operators are polynomials in $\partial_t,\partial_x$ (Section~\ref{sec:theory}) and the ring $\ring$ is commutative, so any two terms commute: $A_1A_2=A_2A_1$. For commuting operators
\[
e^{t(A_1+A_2)}=e^{tA_1}e^{tA_2},\qquad\text{hence}\qquad S=\prod_nS_n;
\]
without commutativity this fails (Section~\ref{sec:noncomm}).

We call
\[
\Delta_n=S_n(t)-\Id
\]
the \emph{increment} produced by term $n$ alone, where $\Id$ is the identity operator and $t$ is fixed; thus $S_n=\Id+\Delta_n$, and $\Delta_nu_0=S_n(t)u_0-u_0$ is exactly the contribution of term $n$ in \eqref{xeq:eint-def}. Products of operators are compositions.
\end{aside}

\begin{notebox}{Fourier variables}
For constant coefficients it is convenient to pass to Fourier variables. With
\[
\hat u(\kappa)=\int u(x)e^{-i\kappa x}\dd x,
\]
where $\kappa\in\R$ is the frequency:
\begin{itemize}[nosep]
\item the derivative $\partial_x$ becomes multiplication by $i\kappa$;
\item each term $A_n$ becomes multiplication by its \emph{symbol} $a_n(\kappa)$ (a polynomial in $i\kappa$);
\item $e^{tA_n}$ becomes multiplication by the number $e^{ta_n(\kappa)}$.
\end{itemize}
Operators turn into numbers and compositions into products of numbers, so for a fixed $\kappa$ every computation below is scalar.
\end{notebox}
\begin{theorem}[{(Interaction term) \led{P1-T2, P1-T3} \proved[]}]\label{xthm:eint}
Assume the $A_n$ commute (e.g.\ constant coefficients). Write $\Delta_n=S_n(t)-\Id$. Then
\begin{equation}\label{eq:eint}
\Eint^{\mathrm{add}}=\Big[\prod_{n=1}^N(\Id+\Delta_n)-\Id-\sum_{n=1}^N\Delta_n\Big]u_0 .
\end{equation}
For $N=2$ this is simply $\Eint=\Delta_1\Delta_2u_0=(S_1-\Id)(S_2-\Id)u_0$. In Fourier variables, with $\hat A_n=a_n(\kappa)$,
\[
\widehat{\Eint}(t,\kappa)=\Big[\prod_n\big(e^{ta_n(\kappa)}\big)-1-\sum_n\big(e^{ta_n(\kappa)}-1\big)\Big]\hat u_0(\kappa).
\]
\end{theorem}
\begin{proof}
Commutation gives $S=e^{t\sum A_n}=\prod_nS_n$. Substitute $S_n=\Id+\Delta_n$ in \eqref{xeq:eint-def}. For $N=2$, expand the product:
\[
\begin{aligned}
S_1S_2&=(\Id+\Delta_1)(\Id+\Delta_2)\\
&=\Id+\Delta_1+\Delta_2+\Delta_1\Delta_2 ,
\end{aligned}
\]
and $\sum_n(S_n-\Id)=\Delta_1+\Delta_2$, so the remainder in \eqref{xeq:eint-def} is $\Delta_1\Delta_2u_0$.

For general $N$, the $\Delta_n$ commute, so
\[
\prod_n(\Id+\Delta_n)=\sum_{S\subseteq\{1,\dots,N\}}\prod_{n\in S}\Delta_n ;
\]
removing the subsets with $|S|=0$ and $|S|=1$ leaves
\[
\Eint=\sum_{|S|\ge2}\Delta_S\,u_0,\qquad \Delta_S=\prod_{n\in S}\Delta_n .
\]
In words: the interaction is the sum of products of at least two increments, and appears only when at least two terms act together.
\end{proof}

\paragraph{A concrete case.}\label{par:advdiff}
Take advection--diffusion, $u_t=-cu_x+Du_{xx}$, with Gaussian datum $u_0=g(x;\sigma^2)$, where $g(x;v)=e^{-x^2/2v}/\sqrt{2\pi v}$. The two flows act as a translation and a Gaussian smoothing:
\[
S_1u_0=g(x-ct;\sigma^2),\qquad S_2u_0=g(x;s^2),\qquad u=g(x-ct;s^2),\qquad s^2=\sigma^2+2Dt .
\]
Hence
\[
\Eint=g(x-ct;s^2)-g(x-ct;\sigma^2)-g(x;s^2)+g(x;\sigma^2).
\]
Figure~\ref{fig:eint}(a) plots it: it is comparable in size to $u$ itself, so it cannot be discarded as a small correction. In Fourier variables the ratio $|\widehat{\Eint}|/|\hat u_0|=|(e^{-ic\kappa t}-1)(e^{-D\kappa^2t}-1)|$ is at most $2$; it vanishes at $\kappa=0$ and whenever $c\kappa t\in2\pi\mathbb Z$ (the translation returns to the identity), see Figure~\ref{fig:eint}(b).

\begin{figure}[!htb]
\centering
\begin{tikzpicture}
\begin{axis}[width=7.6cm,height=5.6cm,xlabel={$x$},xmin=-6,xmax=8,ymin=-0.2,ymax=0.45,
 title={\small (a) $u_t=-cu_x+Du_{xx}$, $t=1$, $c=2$, $D=\tfrac12$},legend style={font=\scriptsize,draw=none,fill=white,fill opacity=0.8,text opacity=1,at={(0.02,0.98)},anchor=north west},
 domain=-6:8,samples=140,no markers]
\addplot[cG,thick,dashed] {gauss(x,1)};
\addplot[cA,thick] {gauss(x-2,2)};
\addplot[cB,very thick] {gauss(x-2,2)-gauss(x-2,1)-gauss(x,2)+gauss(x,1)};
\legend{$u_0$,$u(t)$,$\Eint$}
\end{axis}
\end{tikzpicture}\hfill
\begin{tikzpicture}
\begin{axis}[width=7.0cm,height=5.6cm,xlabel={$\kappa$},xmin=0,xmax=4,ymin=0,ymax=2.1,
 title={\small (b) $|\widehat{\Eint}|/|\hat u_0|$ for several $t$ ($c=2$, $D=\tfrac12$)},legend style={font=\scriptsize,draw=none,at={(0.5,-0.28)},anchor=north,legend columns=3},
 domain=0:4,samples=160,no markers]
\addplot[cA,thick] {sqrt(((cos(deg(2*0.5*x))-1)^2+sin(deg(2*0.5*x))^2)*(exp(-0.5*0.5*x^2)-1)^2)};
\addplot[cC,thick] {sqrt(((cos(deg(2*1.0*x))-1)^2+sin(deg(2*1.0*x))^2)*(exp(-0.5*1.0*x^2)-1)^2)};
\addplot[cB,thick] {sqrt(((cos(deg(2*2.0*x))-1)^2+sin(deg(2*2.0*x))^2)*(exp(-0.5*2.0*x^2)-1)^2)};
\legend{$t=0.5$,$t=1$,$t=2$}
\end{axis}
\end{tikzpicture}
\caption{Advection--diffusion with Gaussian datum ($\sigma^2=1$, $c=2$, $D=\tfrac12$). (a) $u_0$, $u(t)$ and $\Eint$ at $t=1$. (b) The Fourier ratio $|\widehat{\Eint}|/|\hat u_0|$ for three times.}\label{fig:eint}
\end{figure}

\subsection{Cost: $O(N)$ versus $2^N$}
Expanding \eqref{eq:eint} in the $\Delta_n$ produces one monomial per subset of $\{1,\dots,N\}$ with at least two elements: $2^N-N-1$ cross terms (Figure~\ref{fig:lattice}: the red nodes are the cross terms). Writing the \emph{complete table} of cross terms is therefore exponential. Evaluating a \emph{value} is cheap, because we never expand the product. The algorithm uses exactly the product (Fourier) form of Theorem~\ref{xthm:eint},
\[
\widehat{\Eint}(t,\kappa)=\Big[\prod_{n=1}^Ne^{ta_n(\kappa)}-1-\sum_{n=1}^N\big(e^{ta_n(\kappa)}-1\big)\Big]\hat u_0(\kappa),
\]
which for a fixed $\kappa$ needs the $N$ numbers $e^{ta_n(\kappa)}$, $N-1$ products and $N$ subtractions: $O(N)$ operations once $\hat u_0(\kappa)$ is known.
\begin{aside}{Two remarks on cost}
\textbf{What ``value'' means.} Here ``value'' means one Fourier value; recovering $\Eint(t,x)$ at a point of space needs an inverse Fourier transform, whose cost is not counted in the $O(N)$.

\textbf{Where $2^N$ comes from.} The count $2^N-N-1$ refers only to writing out all cross terms one by one; the product form never does so (Step~6 of the algorithm).
\end{aside}
\begin{figure}[!htb]
\centering
\begin{tikzpicture}[x=1.8cm,y=1.1cm,every node/.style={font=\small}]
\node (e) at (0,0) {$\emptyset$};
\node (1) at (-1,1) {$\{1\}$}; \node (2) at (0,1) {$\{2\}$}; \node (3) at (1,1) {$\{3\}$};
\node[text=cB] (12) at (-1,2) {$\{1,2\}$}; \node[text=cB] (13) at (0,2) {$\{1,3\}$}; \node[text=cB] (23) at (1,2) {$\{2,3\}$};
\node[text=cB] (123) at (0,3) {$\{1,2,3\}$};
\foreach \a/\b in {e/1,e/2,e/3,1/12,1/13,2/12,2/23,3/13,3/23,12/123,13/123,23/123} \draw[cG] (\a)--(\b);
\node[anchor=west,align=left,font=\scriptsize] at (1.8,2.2) {\textcolor{cB}{red} = cross terms $\Delta_S=\prod_{n\in S}\Delta_n$\\ $|S|\ge2$: $2^N-N-1$ of them.};
\node[anchor=west,align=left,font=\scriptsize] at (1.8,0.7) {$\Eint=\sum_{|S|\ge2}\Delta_S\,u_0$};
\end{tikzpicture}
\caption{Subset lattice for $N=3$ terms: one node per subset $S$ of terms, one monomial $\Delta_S$ per node with $|S|\ge2$. The count $2^N-N-1$ is what makes the full table exponential.}\label{fig:lattice}
\end{figure}
\begin{proposition}[{(Complexity) \led{P1-T3} \proved[]}]
In the commuting exponential-type class, a value of $\Eint$ costs $O(M(N+\sum d))$ for $M$ evaluation points; aggregated by order $O(MN^2)$; the full table $\Theta(M2^N(L+1)^2)$.
\end{proposition}
\edmark{E1}

